\item Supóngase que un científico tiene razón en pensar que dos cantidades $x$ y $y$ están relacionadas linealmente, esto es $y = mx + b$, por lo menos aproximadamente, para ciertos valores de $m$ y $b$. El científico lleva a cabo un experimento y colecciona datos en forma de puntos $(x_1, y_1), (x_2, y_2), \dots, (x_n, y_n)$, los cuales grafica luego. Estos puntos no se encuentran exactamente en una línea recta, por lo que desea encontrar constantes $m$ y $b$ de manera que la recta $y = mx + b$ ``se ajuste'' a los puntos lo mejor posible. (Véase la figura). Sea $d_i = y_i - (mx_i + b)$ la desviación vertical del punto $(x_i, y_i)$ con respecto a la recta. El método de los mínimos cuadrados determina los valores de $m$ y $b$ de manera que se minimice $\sum_{i=1}^n d_i^2$, la suma de los cuadrados de estas desviaciones. Demuestre que, de acuerdo a este método, la recta de mejor ajuste se obtiene cuando
\[ m \sum_{i=1}^n x_i + bn = \sum_{i=1}^n y_i \]
\[ m \sum_{i=1}^n x_i^2 + b \sum_{i=1}^n x_i = \sum_{i=1}^n x_i y_i \]
De manera que la recta se encuentra resolviendo estas dos ecuaciones en las dos incógnitas $m$ y $b$.
